
\section{مقایسه مدل‌های \lr{Hodgkin-Huxley} و \lr{Morris-Lecar}}

در این بخش، رفتار دینامیکی مدل کامل \lr{Hodgkin-Huxley} (\lr{HH}) و مدل \lr{Morris-Lecar} (\lr{ML}) مقایسه می‌شود. مدل \lr{HH} شامل چهار متغیر حالت $V$، $m$، $h$ و $n$ است. برای بررسی صفحه فاز مدل \lr{ML}، از نسخه کاهش‌یافته دوبعدی آن استفاده شده است. در این تقریب، فعال‌سازی کانال کلسیم سریع در نظر گرفته می‌شود؛ یعنی $m=m_\infty(V)$ و تنها متغیرهای حالت، ولتاژ غشا $V$ و متغیر بازیابی $w$ هستند.

در شبیه‌سازی‌های مدل \lr{HH}، جریان خارجی برابر
 $I_{ext}=10$ 
   و در مدل \lr{ML} برابر $I_{ext}=90$  در نظر گرفته شده است (این مقادیر با مشورت با هوش مصنوعی انتخاب شدند). 

\subsection{مقایسه  فاز}

شکل \ref{fig:phase-portrait} مسیر حالت مدل \lr{HH} را در صفحه $V$-$n$ و مسیر مدل کاهش‌یافته \lr{ML} را در صفحه $V$-$w$ نشان می‌دهد. در مدل \lr{HH}، مسیر پس از گذر از یک دوره کامل به نزدیکی حالت قبلی بازمی‌گردد و یک حلقه بسته تشکیل می‌دهد. این به این معنی است که نورون پس از هر اسپایک وارد مرحله بازیابی می‌شود و سپس دوباره اسپایک می‌زند. 

در مقابل، مسیر مدل \lr{ML} به‌صورت مارپیچی به سمت یک نقطه مشخص حرکت می‌کند. این الگو نشان می‌دهد که دستگاه دارای یک نقطه تعادل پایدار است. بنابراین، در شرایط شبیه‌سازی حاضر، نورون پس از یک پاسخ گذرا یا یک اسپایک اولیه، به حالت تعادل بازمی‌گردد و بصورت متناوب اسپایک نمی‌زند.

\begin{figure}[H]
	\centering
	\includegraphics[width=\textwidth]{pics/1_phase_portrait.png}
	\caption{
		مقایسه  فاز مدل‌ها: مسیر مدل \lr{Hodgkin-Huxley} در صفحه ولتاژ-$n$ (چپ) و مسیر مدل کاهش‌یافته \lr{Morris-Lecar} در صفحه ولتاژ-$w$ (راست).}
	\label{fig:phase-portrait}
\end{figure}

البته نمودار $V$-$n$ برای مدل \lr{HH} تنها یک تصویر دوبعدی از دستگاه چهاربعدی اصلی است. برای تحلیل کامل‌تر، باید تغییرات متغیرهای $m$ و $h$ نیز در نظر گرفته شوند، اما چون هدف نمایش نمودار بود، تنها در دو بعد این نمودار نشان داده شد و همین نمودار نیز شباهت و تفاوت دو مدل را به خوبی نمایش می‌دهد

\subsection{بررسی نال‌کلاین‌ها}

برای مدل کاهش‌یافته \lr{Morris-Lecar}، نال‌کلاین‌ها از قرار دادن مشتق‌ها برابر صفر به‌دست می‌آیند:
\begin{align*}
	w_V(V) &= \frac{I_{ext}-g_{Ca}m_\infty(V)(V-V_{Ca})-g_L(V-V_L)}
	{g_K(V-V_K)},\\
	w_w(V) &= w_\infty(V).
\end{align*}

تقاطع این دو منحنی، نقطه یا نقاط تعادل دستگاه را تعیین می‌کند. همان‌طور که در شکل \ref{fig:ml-nullclines} دیده می‌شود، نال‌کلاین ولتاژ و نال‌کلاین متغیر $w$ تنها در یک نقطه یکدیگر را قطع می‌کنند. جهت میدان برداری در اطراف این نقطه نشان می‌دهد که مسیرها، حتی پس از دور شدن از تعادل، تمایل دارند دوباره به آن نزدیک شوند. بنابراین این نقطه یک تعادل پایدار از نوع کانون مارپیچی است.

\begin{figure}[H]
	\centering
	\includegraphics[width=0.72\textwidth]{pics/2_nullclines.png}
	\caption{نال‌کلاین‌های ولتاژ و متغیر $w$ در مدل کاهش‌یافته \lr{Morris-Lecar}. پیکان‌ها علامت و جهت $(\dot V,\dot w)$ را در نواحی مختلف نشان می‌دهند و محل تقاطع دو منحنی، نقطه تعادل دستگاه است.}
	\label{fig:ml-nullclines}
\end{figure}

برای مقایسه با \lr{HH}، در شکل \ref{fig:hh-nullclines} نال‌کلاین‌های یک تقریب دوبعدی از این مدل رسم شده‌اند. این تقریب فقط با هدف نمایش صفحه فاز انجام شده است و در آن فرض می‌شود
\[
m=m_\infty(V), \qquad h=1-n.
\]
با این فرض، فقط متغیرهای ب $V$ و $n$  باقی می‌مانند. و برای بدست آوردن نال کلاین‌ها باید مشتق را برابر صفر کرد. 

\begin{figure}[H]
	\centering
	\includegraphics[width=0.72\textwidth]{pics/2_hh_nullclines.png}
	\caption{نال‌کلاین‌های مدل ساده‌شده \lr{Hodgkin-Huxley} با فرض $m=m_\infty(V)$ و $h=1-n$. پیکان‌ها جهت $(\dot V,\dot n)$ را نشان می‌دهند.}
	\label{fig:hh-nullclines}
\end{figure}

در مدل کامل \lr{HH}، نقطه تعادل باید به‌طور همزمان شرایط زیر را ارضا کند:
\begin{align*}
	\dot V&=0, &m&=m_\infty(V), &h&=h_\infty(V), &n&=n_\infty(V).
\end{align*}
در نتیجه، نال‌کلاین‌های مدل کامل در فضای چهاربعدی تعریف می‌شوند و نمی‌توان آن‌ها را مستقیماً مانند مدل دوبعدی \lr{ML} در یک نمودار صفحه‌ای نمایش داد. به همین دلیل، از مدل ساده شده 

\lr{HH}
استفاده شد تا بتوان نمودار آن را رسم کرد و با 
مدل 
\lr{ML}
مقایسه کرد.


\subsection{نقاط تعادل و پایداری}

حل عددی تقاطع نال‌کلاین‌های مدل \lr{ML}، یک نقطه تعادل تقریبی در مکان زیر به‌دست می‌دهد:
\begin{center}
	$V^*\simeq -46.3\ \mathrm{mV}$ و $w^*\simeq 0.22$.
\end{center}

ژاکوبین دستگاه کاهش‌یافته در این نقطه تقریباً برابر است با
\[
J \simeq
\begin{pmatrix}
	-3.30 & -301.6\\
	0.0677 & -1.97
\end{pmatrix}.
\]
مقادیر ویژه این ماتریس نیز برابرند با
\[
\lambda_{1,2}\simeq -2.64 \pm 4.46i.
\]
منفی بودن قسمت حقیقی مقادیر ویژه نشان می‌دهد که نقطه تعادل پایدار است. از آن‌جا که مقادیر ویژه مختلط‌اند، بازگشت سیستم به تعادل به‌صورت نوسانی انجام می‌شود. به همین دلیل، مسیر مدل \lr{ML} در صفحه فاز به شکل مارپیچی به نقطه تعادل نزدیک می‌شود. 

برای مدل ساده‌شده \lr{HH} نیز ژاکوبین عددی دستگاه دوبعدی $(V,n)$ در محل تقاطع نال‌کلاین‌ها محاسبه شد. سه نقطه تعادل زیر به‌دست آمد:
\begin{align*}
	(V_1,n_1)&\simeq(-58.17,0.426), &
	J_1&\simeq\begin{pmatrix}1.95&-228.26\\0.00318&-0.200\end{pmatrix},\\
	(V_2,n_2)&\simeq(-50.99,0.536), &
	J_2&\simeq\begin{pmatrix}14.17&-726.09\\0.00333&-0.226\end{pmatrix},\\
	(V_3,n_3)&\simeq(-21.32,0.828), &
	J_3&\simeq\begin{pmatrix}4.96&-10037.92\\0.00232&-0.421\end{pmatrix}.
\end{align*}

مقادیر ویژه متناظر، به‌ترتیب، تقریباً برابر $1.54,0.22$، $14.00,-0.06$ و $2.27\pm4.00i$ هستند. بنابراین، نقطه اول ناپایدار، نقطه دوم از نوع زینی و نقطه سوم یک کانون مارپیچی ناپایدار است. پس در تقریب دوبعدی مورد استفاده، هیچ‌یک از نقاط تعادل پایدار نیستند.

این نتیجه با رفتار تناوبی مشاهده‌شده در شکل \ref{fig:phase-portrait} سازگار است؛ زیرا مسیرها به‌جای همگرا شدن به یکی از نقاط تعادل، از نواحی اطراف آن‌ها دور شده و به سمت چرخه تناوبی حرکت می‌کنند. 

\subsection{خطی‌سازی و تقریب}

برای مدل کاهش‌یافته \lr{ML}، اگر
\[
\delta x=(\delta V,\delta w)^T
\]
انحراف کوچک از نقطه تعادل
\[
x^*=(V^*,w^*)^T
\]
باشد، دینامیک خطی‌شده در همسایگی تعادل به‌شکل زیر نوشته می‌شود:
\[
\dot{\delta x}=J\,\delta x.
\]

شکل \ref{fig:ml-linearization} پاسخ مدل غیرخطی و مدل خطی‌شده را برای فاصله اولیه
\[
\delta x=(0.1,0.01)^T
\]
مقایسه می‌کند. در ابتدای حرکت، دو پاسخ تقریباً بر هم منطبق هستند؛ چون خطی‌سازی تنها در نزدیکی نقطه تعادل معتبر است. با افزایش فاصله از تعادل، اثر غیرخطی بودن جریان‌های یونی آشکار می‌شود و اختلاف کوچکی میان دو مسیر به وجود می‌آید. با این وجود، هر دو پاسخ در نهایت کاهش یافته و به سمت نقطه تعادل حرکت می‌کنند.

\begin{figure}[H]
	\centering
	\includegraphics[width=\textwidth]{pics/4_linearization.png}
	\caption{مقایسه پاسخ غیرخطی و خطی‌شده مدل \lr{Morris-Lecar}: ولتاژ برحسب زمان (چپ) و مسیر در صفحه فاز (راست). فاصله اولیه از تعادل $\delta x=(0.1,0.01)^T$ است.}
	\label{fig:ml-linearization}
\end{figure}

برای مدل ساده‌شده \lr{HH}، خطی‌سازی در همسایگی نقطه سوم، یعنی $(V_3,n_3)$، انجام شده است. فاصله اولیه در این حالت برابر
\[
\delta x=(0.1,0.001)^T
\]
در نظر گرفته شد. همان‌طور که در شکل \ref{fig:hh-linearization} دیده می‌شود، پاسخ خطی و غیرخطی در لحظات ابتدایی به هم نزدیک‌اند، اما با گذشت زمان از یکدیگر فاصله می‌گیرند. علت اصلی این موضوع، ناپایداری نقطه تعادل است؛ زیرا قسمت حقیقی مقادیر ویژه، یعنی $2.27 \pm 4.00i$، مثبت است و انحراف‌های کوچک را تقویت می‌کند.

در مقایسه با \lr{HH}، مدل \lr{ML} در اطراف نقطه تعادل خود تقریب خطی مناسب‌تری ارائه می‌دهد. وجود تعادل پایدار باعث می‌شود خطاهای کوچک در تقریب خطی رشد نکنند و مسیر به سمت ناحیه تعادل باقی بماند. در مدل ساده‌شده \lr{HH}، حتی با وجود فاصله اولیه کوچک‌تر، ناپایداری تعادل و غیرخطی بودن شدیدتر دینامیک باعث می‌شود تقریب خطی در زمان کوتاه‌تری از پاسخ اصلی فاصله بگیرد.

\begin{figure}[H]
	\centering
	\includegraphics[width=\textwidth]{pics/5_hh_linearization.png}
	\caption{مقایسه مدل غیرخطی و خطی‌شده \lr{Hodgkin-Huxley} ساده‌شده در همسایگی نقطه ناپایدار $(V_3,n_3)$ با فاصله اولیه $\delta x=(0.1,0.001)^T$. دو پاسخ در ابتدا نزدیک‌اند، اما با گذر زمان واگرا می‌شوند.}
	\label{fig:hh-linearization}
\end{figure}
